\appendix
et comparer des cas du benchmark
\section*{Metaporous response obtained using Comsol}
\mm{enlever ça}

In this appendix, the Comsol model used for comparison with the proposed method is described. For the considered cases, the geometry is reproduced within the software and the same material parameters are used. The surrounding medium at both ends of the structure is modeled as an air layer of width $d$ and height $4H$ terminated at each side by a Perfectly Matched Layer (PML) with height $4H$. The geometry uses triangular elements for the geometry and mapped square elements for the PML, with characteristic sizes for the mesh $l_{car} = c_0/(15f_{max})$. The physical interfaces are modeled using the stock physics libraries from Comsol for poroelasticity, pressure acoustics and solid mechanics, respectively for PEM, fluid and elastic layers. Each are coupled with their appropriate Multiphysics couplings. 

The excitation is imposed by a background pressure field and periodic conditions are applied for each layer at the left and right edges. The reflection and transmission coefficients of the structure are evaluated using an average operators, defined along a line above and below the structure in each surrounding fluid layer, respectively $l_1$ (above the structure, in the upper fluid domain) and $l_2$ (below the structure, in the lower fluid structure). This allows to use the pressure fields to write,
\begin{equation}
    R = e^{ik_y l_1} \frac 1 D \int_0^D {p_s(x, y=l_1)} e^{-ik_x x} dx, \quad
    T = e^{ik_y l_2} \frac 1 D \int_0^D {p_t(x, y=l_2)} e^{-ik_x x} dx,
\end{equation}
% \begin{equation}
 %   R_q=\frac{1}{|\mathcal{L}|}\mathrm{e}^{-\mathrm{i}k_2^{(q)}x_2(\mathcal{L})}\int_{\mathcal{L}}p_\mathrm{s}(\mathbf{x},\omega)\mathrm{e}^{-\mathrm{i}k_1^{(q)}x_1}\mathrm{d}x_1,
% \end{equation}
where the scattered and total pressure fields are denoted as $p_s$ and $p_t$, respectively. The absorption coefficient is finally calculated as $\alpha = 1 - |R|^2 - |T|^2$.